% ECONOMICS_PROBLEM: {"id": "L41-605", "title": "Prospective merger prediction versus ex-post effects", "jel_code": "L41", "source_id": "OP-0071"}
% !TeX program = xelatex
\documentclass[11pt,a4paper]{article}
\usepackage[margin=23mm]{geometry}
\usepackage{fontspec}
\setmainfont{Latin Modern Roman}
\usepackage{amsmath,amssymb,mathtools}
\usepackage{xcolor,enumitem,booktabs,longtable,array,xurl,hyperref}
\definecolor{muted}{HTML}{53636C}
\hypersetup{colorlinks=true,urlcolor=blue,linkcolor=blue}
\setlength{\parindent}{0pt}
\setlength{\parskip}{5pt}
\newcommand{\R}{\mathbb R}
\newcommand{\E}{\mathbb E}
\newcommand{\N}{\mathbb N}
\newcommand{\ind}{\mathbf 1}
\newcommand{\Var}{\operatorname{Var}}
\newcommand{\Cov}{\operatorname{Cov}}
\newcommand{\supp}{\operatorname{supp}}
\DeclareMathOperator*{\argmin}{arg\,min}
\DeclareMathOperator*{\argmax}{arg\,max}
\newcommand{\entrymeta}[3]{{\sffamily\small\color{muted}\textbf{\texttt{#1}}\quad|\quad #2\par #3\par}}
\newcommand{\Problemref}[1]{\texttt{#1}}
\newcommand{\JELref}[2]{\texttt{#2} (JEL \texttt{#1})}
\begin{document}
\section*{Prospective merger prediction versus ex-post effects}
\textbf{Problem ID:} \texttt{L41-605}\quad \textbf{JEL:} \texttt{L41}\par
% BEGIN ECONOMICS STATEMENT
\label{problem:OP-0071}
\label{atlas:prob:I1}

\noindent\textbf{Original identifier:} I1 (atlas item 71). \textbf{Classification:} Prospective causal prediction research program. \textbf{Original tractability label:} Medium (editorial, not a complexity result).

\subsection*{Definitions and assumptions}

For proposed merger $j$, define a preregistration date $t_j$, baseline information $\mathcal I_j$, and a policy $M_j\in\{0,1\}$ indicating consummation under a specified remedy regime. Let $Y_{jh}(m)$ be the market outcome $h$ periods later under policy $m$, with components for quality-adjusted prices, quantities, entry, and consumer surplus. The counterfactual specifies common macroeconomic shocks and the conduct, entry, and efficiency responses allowed under each policy. A prospective model delivers a predictive distribution $Q_j$ for the causal effect $\Delta_{jh}=Y_{jh}(1)-Y_{jh}(0)$ using only information available at $t_j$.

\subsection*{Formal research question}

Construct and validate a protocol evaluating $Q_j$ against retrospectively identified $\Delta_{jh}$, including uncertainty in that retrospective identification. For a declared loss function $L$, target
\[
R_h=\mathbb E_{j\sim F_*}[L(Q_j,\Delta_{jh})],
\]
where $F_*$ is the specified population of proposed mergers. Characterize sufficient sampling, causal-identification, and equilibrium restrictions for identification of $R_h$. Test calibration of prespecified prediction intervals and compare models on held-out merger cohorts. If only selected consummated cases are observed, derive bounds or transport corrections for the target proposal population.

\subsection*{Interpretation and scope}

A before--after price change is not the merger's causal effect: costs, demand, and rival behavior change contemporaneously. Even a well-designed retrospective may identify effects only for consummated mergers, whereas a prospective benchmark targets cases including blocked or abandoned transactions. This selection issue requires a defined target and assumptions, not a mechanical comparison with realized prices. Forecasts should preserve their original information sets and confidence statements, avoiding revisions after the event. Static merger simulation and dynamic entry models may answer different horizon questions. Quality and product removal require welfare-consistent measurement rather than comparison of prices for a changing set of products. A benchmark can score partial predictions when the retrospective delivers an identified set, using a prespecified rule that does not reward uninformative intervals. Predictive success for prices alone does not establish correct entry or innovation predictions. Model performance is an empirical quantity, not a universal yes-or-no property of economics.

\subsection*{Source audit and status}

Nevo (2000) is contextually resolved to the cereal-merger paper [S1], because that author-year is not unique. [S2] directly compares simulated and observed airline prices; [S3] provides consummated-merger evidence. Original link [40], [S4], confirms the agency function and has no role as a scholarly open-problem source. The prediction-risk target is editorial and includes selection and retrospective uncertainty omitted from the original short question.

\subsection*{References}

\begin{enumerate}[label={[S\arabic*]},leftmargin=*]

\item Aviv Nevo (2000). \emph{Mergers with Differentiated Products: The Case of the Ready-to-Eat Cereal Industry}. RAND Journal of Economics 31(3), 395--421. \url{https://escholarship.org/uc/item/1d53t6ts}\par \textit{Source role:} Merger simulation example; author-year selection is contextual. \textit{Locator:} Abstract and article metadata.

\item Craig Peters (2006). \emph{Evaluating the Performance of Merger Simulation: Evidence from the U.S. Airline Industry}. Journal of Law and Economics 49(2), 627--649. \url{https://doi.org/10.1086/505369}\par \textit{Source role:} Comparison of merger predictions and realized prices. \textit{Locator:} Abstract and article metadata.

\item Orley Ashenfelter, Daniel Hosken, and Matthew Weinberg (2014). \emph{Did Robert Bork Understate the Competitive Impact of Mergers? Evidence from Consummated Mergers}. Journal of Law and Economics 57(S3), S67--S100. \url{https://doi.org/10.1086/675862}\par \textit{Source role:} Merger retrospective evidence. \textit{Locator:} Abstract and article metadata.

\item Federal Trade Commission (accessed 2026-10-08). \emph{Merger Review}. Agency topic page. \url{https://www.ftc.gov/news-events/topics/competition-enforcement/merger-review}\par \textit{Source role:} Institutional background; not proof of a research frontier. \textit{Locator:} Abstract and article metadata.

\end{enumerate}

\subsection*{Common conventions: Source mathematical conventions}

Unless an entry states otherwise, agent and firm sets are finite. State and action spaces are measurable, strategies and outcome maps are measurable, probability laws are specified, and expectations exist for the targets under discussion. Infinite-horizon discount factors lie in $(0,1)$ and discounted objectives are finite. Norms, tolerances and complexity input sizes must be specified for computational comparisons. Constants or primitives that appear locally are inputs, not empirical facts asserted by this catalogue.

\textbf{Identification.} A statistical model $m\in\mathcal M$ implies an observable law $P_m$ and target $\theta(m)$. At law $P$, its identified set is
\[
\Theta(P;\mathcal M)=\{\theta(m):m\in\mathcal M,\ P_m=P\}.
\]
Point identification holds when this set is a singleton, or equivalently when $P_{m_1}=P_{m_2}$ implies $\theta(m_1)=\theta(m_2)$ for all compatible models. Sharp bounds mean bounds attained, or approached when the boundary is not attained, by compatible models. Writing this set is a definition, not a solution: the substantive task is to establish informative restrictions and derive or estimate the set. An unrestricted-confounding model may give only trivial bounds.

\textbf{Inference.} Identification, estimability and valid uncertainty quantification are separate questions. Fix $\alpha\in(0,1)$. A confidence procedure $C_n$ over a declared law class $\mathcal P$ has asymptotic uniform coverage at level $1-\alpha$ when
\[
\liminf_{n\to\infty}\inf_{P\in\mathcal P}\Pr_P\{\theta(P)\in C_n\}\geq1-\alpha.
\]
For set-identified targets the coverage event must be defined separately, for example coverage of the full identified set. If an entry uses identification-indexed models instead of uniquely indexed laws, the same coverage convention is applied over those models. Pointwise limits do not establish uniform validity, and asymptotic validity does not guarantee finite-sample precision. Sample sizes, dependence structures and weak-identification sequences are part of each application.

\textbf{Counterfactuals.} $Y_i(a)$ is a potential outcome under a specified intervention, including its timing, implementation and versions. It is not $Y_i$ conditional on voluntarily choosing $a$. A full intervention vector permits interference; shorthand scalar treatment notation does not silently impose no interference. Assignment, selection, attrition, anticipation and measurement must be specified. A claim about an instrument's local effect is not automatically a population policy effect.

\textbf{Economic equilibrium.} A counterfactual model includes best responses, constraints, feasibility, market clearing, state transitions and expectations consistent with the stated information. When several equilibria exist, either a declared selector is an input or the target is set-valued. Comparisons across policy regimes must use the stated selection convention. Reduced-form relationships are not presumed invariant after an equilibrium-changing intervention.

\textbf{Optimization and welfare.} Optimization is over the stated feasible set. If attainment is not established, $\sup$ or $\inf$ and $\varepsilon$-optimal choices replace an asserted maximizer or minimizer. For example, with value $V=\sup_{a\in\mathcal A}W(a)<\infty$, an $\varepsilon$-optimal set is $\{a\in\mathcal A:W(a)\geq V-\varepsilon\}$ for $\varepsilon>0$. Benefits, costs, risk and health or environmental outcomes can be aggregated only using declared units and conversion weights. Interpersonal utility weights, the treatment of fiscal transfers and foreign profits, discounting, and the behavioral welfare criterion are explicit inputs. A comparative-static derivative additionally requires existence and appropriate differentiability of the selected counterfactual.

\textbf{Sources.} Local labels [S1], [S2], and so forth restart in every question. Locators identify the relevant abstract, model, section or result at the level actually checked; a bibliographic or abstract-level verification is not represented as inspection of an entire proof. Working papers are distinguished from later publications and changing versions. Source summaries are paraphrased. All URL checks and editorial formulations refer to the audit date stated above.
% END ECONOMICS STATEMENT
\end{document}
